% Supplementary text for an anonymous submission. Repository URL, author names,
% original Git commit and private workspace paths must remain in an author-only
% provenance ledger, not in this review manuscript.
\section{Public-release construction audit and reproducibility}
\label{supp:public-audit}

We used the unchanged VDJdb release dated 1 June 2023. The initial table contains
92,771 records. Sequentially requiring human origin, the T-cell receptor beta
chain, major histocompatibility complex class I, exact two-field HLA-A/B/C
annotation syntax, a canonical amino-acid CDR3 beginning with C and ending with F or W,
and a canonical peptide of length 8--15 retained 85,329, 49,686, 47,005, 35,278,
35,242, and 35,237 records, respectively. HLA resolution was not imputed. We did
not silently impose a CDR3 length threshold. After confidence filtering and
deduplication of exact beta/peptide/HLA triples, the score-at-least-one primary
cohort contained 3,182 triples, 3,061 distinct beta sequences, 253 peptides, and
41 source HLA labels. The primary analysis preserves source annotations and
does not assume that syntactic validity establishes current nomenclature
validity. The score-at-least-zero and score-three sensitivity cohorts
contained 29,306 and 320 triples, respectively. Their pre-deduplication counts
were 35,237 and 594; the primary pre-deduplication count was 4,538.

The known-positive reference included all 29,306 eligible triples at score zero
or above, regardless of the narrower analysis cohort. We separately counted
exact triple collisions and beta/peptide collisions without considering HLA.
These are contradictions with a recorded positive reference, not estimates of
the complete biological false-negative rate. A second sensitivity analysis
restricted the reference to confidence scores of at least one. Database
cross-reactivity labels and confidence scores are taken as provided, without
experimental revalidation.

\subsection{Execution of source logic}
The source sampler's class body was extracted from its abstract syntax tree and
executed unchanged. A minimal numerical interface replaced only tensor
stacking, concatenation, and conversion; unique identity tags replaced sequence
embeddings. The selected antigen identities were recovered from these tags.
No protein language model, neural checkpoint, or original study embedding was
executed. This design audits the actual sequence-selection and label-generation
logic without interpreting synthetic tags as learned representations. An
anonymous source snapshot and file hashes support reproduction; identifying
repository metadata are withheld from the anonymous manuscript.

For each cohort, we fixed the global fallback pool to the first occurrence of
each distinct peptide/HLA pair in source order. We independently varied loader
order among release order, peptide/HLA-sorted order, and random row order.
Source-order ties in sorted loading were preserved. The loader permutation used
NumPy's default random generator with seed $100000+s$, and the source sampler
used Python's random generator with seed $s$, for $s=0,\ldots,29$. All cohorts
used batch size 64; the primary cohort additionally used batch sizes 16 and 256.
Changing loader order did not reorder the fallback pool. The resulting 540
stochastic runs, together with one deterministic corrected construction,
generated 3,770,361 negative records in the main analysis. Repeated records
across runs are Monte Carlo replicates, not additional biological observations.

\subsection{Marginal diagnostics}
For equally sized positive and generated-negative classes, marginal
total-variation distance was computed as
\begin{equation}
  \operatorname{TV}(P_+,P_-)=\frac12\sum_c\lvert P_+(c)-P_-(c)\rvert.
\end{equation}
Categories $c$ were peptide identity, HLA identity, joint peptide/HLA identity,
or beta-chain identity. The identity-oracle diagnostic assigned category $c$
the value $n_+(c)/(n_+(c)+n_-(c))$ and evaluated AUROC on the same constructed
dataset, with half credit for tied scores. It is deliberately a resubstitution
measure of marginal class separability. It does not measure held-out prediction
or generalization to unseen peptides or receptors. We report means, sample
standard deviations, ranges, and empirical 2.5th and 97.5th percentiles across
construction seeds; these are not population confidence intervals.

\subsection{Maximum retained closed assignment}
Within an HLA stratum with $N$ positive rows, let $C_{ij}=0$ when assigning the
antigen of donor row $j$ to receptor row $i$ creates a triple absent from the
known-positive reference. Set $C_{ii}=1$ to encode abstention, and assign cost
$N+1$ to every other forbidden edge. A minimum-cost complete assignment is
computed using SciPy's \texttt{linear\_sum\_assignment}. The identity assignment
has cost $N$, so no optimum can contain an edge of cost $N+1$. Its non-diagonal
assignments form disjoint cycles. Removing diagonal self-loops therefore retains
the identical source and donor row set. Minimizing the cost minimizes the number
of abstentions, yielding maximum retained cardinality among closed within-HLA
permutations obeying the positive-exclusion rule. This is an application of the
standard assignment problem, not a new optimization algorithm.

We verified zero reference-positive collisions and exact equality of beta,
peptide, HLA, and joint peptide/HLA class marginals in the returned construction.
The primary solution retained 2,691 of 3,182 positive rows (84.57\%), 236 of 253
peptides, and 22 of 41 HLA alleles. Retention is selective: single-peptide HLA
strata cannot supply a different peptide under the stated constraints, and
cross-reactivity can impose additional restrictions. Original and corrected
samplers were therefore compared both descriptively on the original population
and directly on the same retained positive rows. This construction does not
certify the remaining generated pairs as experimentally verified nonbinders.

\subsection{HLA annotation sensitivity}
An independent nomenclature review identified two source labels requiring
additional scrutiny: \texttt{HLA-A*08:01} was unresolved in the current
nomenclature comparison, and \texttt{HLA-A*24:01} corresponds to a deleted legacy
annotation. The official deletion record reports a sequence error in March
1995. We did not infer biological binding consequences from nomenclature status
and did not replace either label with another allele. A secondary analysis
removed both exact labels from the analysis cohorts and the broad known-positive
reference, leaving 3,108 primary positive triples. We reran all three loader
orders at batch size 64 for 30 seeds, on both the full clean cohort and the
common retained cohort.

The clean full-cohort mean known-positive collision rates were 2.8882\%,
0.9331\%, and 0.2445\% for release, peptide-sorted, and shuffled orders,
respectively. Corresponding peptide total-variation distances were 0.182937,
0.439168, and zero; identity-oracle AUROCs were 0.642295, 0.793573, and 0.500000.
The exact corrected construction retained 2,665 of 3,108 rows (85.75\%), with
zero known-positive collisions and exact equality of all four audited marginal
distributions. These results preserve the qualitative primary conclusions.
Deleting records changes batch boundaries, so differences from the raw-source
audit are not interpreted as a biological effect of those HLA annotations.

\subsection{Strict peptide-component sensitivity}
The 253 primary-cohort peptides formed 221 connected components under unit-cost
Levenshtein edit distance at most two. For five prespecified seeds, components
were permuted with NumPy seed $200000+s$, and 132/44/45 components were assigned
to training, validation, and test partitions. This approximates 60/20/20 by
component count; record counts are reported separately. Exhaustive comparison
verified a minimum cross-partition edit distance of three for every partition.
Test populations contained 1,022, 877, 643, 372, and 593 rows. The corrected
within-test constructions retained 509, 546, 380, 271, and 190 rows, respectively
(49.80\%, 62.26\%, 59.10\%, 72.85\%, and 32.04\%). Each original test
construction was audited under three loader orders and ten sampler seeds on
both full and retained test populations. No model was trained on these
partitions, so this sensitivity experiment tests construction invariants under
peptide separation and makes no out-of-distribution performance claim.

\subsection{Frozen manifests and release restrictions}
Frozen manifests identify the zero-based source row for each receptor and the
source row providing each negative peptide/HLA identity. Row zero denotes the
first data row after the TSV header. An antigen source row may identify the
first equivalent peptide/HLA pair rather than an indistinguishable donor object
inside a batch. A sentinel index of $-1$ encodes the original source sampler's
hardcoded fallback. We independently reconstructed every frozen negative triple
from the unchanged source table and verified exact agreement with the local
analysis table. Sequence-derived materializations are regenerated locally;
redistribution contains source-row manifests and aggregate measurements, with
the unchanged original database archive and its license retained separately.
